\section{Reviewer 8}

\subsection*{R8-1 --- First-round items}

\begin{rquote}
Most of my first-round comments have been addressed adequately. \ldots
\end{rquote}

\resp Thank you for the detailed check of each item. We have not changed anything the
first-round response established; the revision below is limited to the seven points of this
round.

\subsection*{R8-2 --- Section 4 has grown; state the limitations in the abstract}

\begin{rquote}
Section 4 is reframed as design motivation as requested (comment 7), but it has become longer
rather than shorter \ldots I accept the explanations for comments 9 and 10, but the
limitations should also be stated briefly in the abstract.
\end{rquote}

\paragraph{Section 4 length.} You are right that the reframing added length. Counting the
section's own source with the LaTeX markup stripped (\texttt{detex}, comments removed):
2{,}679 words at first submission, 3{,}165 in the version you reviewed, \textbf{2{,}724} now.
That is $-441$ words ($-13.9\%$) against the reviewed version. It is still $+45$ words
($+1.7\%$) against the first submission, and those 45 are not padding: they are the
cross-term correction of R3-1 (the expansion is now an identity with the first-order term
displayed, and the drift-ratio assumption stated rather than absorbed into a remainder) and
the bound-slack audit of R7-W6, both of which you accepted in the previous round. Removing
either would reopen an accepted correction.

The length was recovered by compression, not by deleting content: the duplicated
design-motivation caveat was removed (it now appears once, where the design rule is stated),
the proof roadmap and the stationarity remark were tightened, and the survey of prior
forgetting bounds --- Nguyen, Evron, Karakida, Doan, OGD --- moved in full, with its
citations, to a new Appendix~C, where Sec.~4 now points in one sentence.
We did not cut the cross-term discussion or the slack audit.

\paragraph{Limitations in the abstract.} The abstract now ends with the limitations as three
clauses rather than a gesture: \emph{one backbone and one adapter site}, \emph{LoRA on the
attention projections untested}, and \emph{training optimizes a linear head rather than the
NCM rule used to evaluate}. The memory trade-off was already in the abstract as a cost of the
design; we kept it separate from these three so that the scope limits are not read as a cost
statement. Adding these clauses, and the attribution of R8-3, took the abstract past the
journal's 250-word limit, so the rest of it was tightened to fit --- no claim was dropped from
the paper and no number changed.

\change Sec.~4 (compressed, $-441$ words); new Appendix~C; abstract
(limitations sentence).

\subsection*{R8-3 --- The margin comes from regularizing a shared adapter, not from the Fisher weighting}

\begin{rquote}
FOLoRA now shows no advantage over tuned diagonal EWC-LoRA (including at T=50), no ceiling
advantage over a tuned equal-weight penalty, and ties a tuned isotropic displacement penalty
(Exp.~A4). The remaining margin over the parameter-efficient family therefore appears to come
from regularizing a shared adapter rather than from the Fisher weighting. Please consider
reflecting this in the title, abstract and highlights, and reporting the isotropic penalty
alongside FOLoRA in the comparison with the parameter-efficient family.
\end{rquote}

\resp We agree with the reading, including the part that costs us the mechanism claim, and the
paper is now written that way. Three changes.

\paragraph{Title.} The title now names the object the evidence supports --- one shared adapter
--- rather than leaning on the Fisher adjective while the paper reports that the weighting is
not what produces the margin. It reads: \emph{``FOLoRA: Fisher-Orthogonalized Low-Rank
Adaptation of One Shared Adapter for Continual Learning''}. The acronym is retained because it
is the identifier used throughout the paper, the code and the repository; the phrase after the
colon now states what is being adapted.

\paragraph{Abstract.} The results sentence was rewritten to attribute the margin correctly.
Where it previously said ``its positive result is against the parameter-efficient family, where
it improves significantly over L2P, CODA-Prompt and InfLoRA \ldots'', it now states that
FOLoRA shows no significant difference from tuned EWC-LoRA or from a tuned isotropic
displacement penalty of matched RMS, and that the margin over the per-task-module family is
therefore attributable to regularizing a \emph{shared} adapter rather than to the Fisher
weighting specifically. The size of the tie is stated in the abstract: ``nor from a
direction-agnostic isotropic penalty of matched RMS (within $0.1$ point on both)''.

\paragraph{Highlights.} The fifth highlight was rewritten from a bare win-list to the
attribution: \emph{``The margin comes from sharing the adapter, not Fisher weighting;
isotropic ties.''} The second was narrowed to the comparison it comes from --- eigenvalue
weighting against the uniform-weighting arm on the same grid, where the $3$--$10\times$
difference in operating strength is measured --- because the general form (``Fisher weighting
sets the operating strength'') would now sit awkwardly beside the isotropic tie. Both stay
inside the journal's 85-character limit per highlight, which is why they are phrased
compactly. The first, third and fourth are unchanged.

\paragraph{The isotropic penalty in the comparison.} As in R1-1, it is now a row of
Table~1, immediately above FOLoRA, and is discussed in
Sec.~5.2 alongside the other baselines. Its numbers are
$78.27\pm0.38$ (CIFAR-100, $\lambda{=}3$) and $66.89\pm0.29$ (ImageNet-R, $\lambda{=}3$)
against FOLoRA's $78.37\pm0.84$ and $66.89\pm0.41$ on the same three seeds; the paired
differences are inside the noise on both metrics ($\Delta{=}+0.00$, $p{=}0.98$ on ImageNet-R
accuracy), and on the gap metric the point estimate is nominally in its favour on CIFAR-100
($\Delta{=}-0.22$, $p{=}0.49$) and nominally against it on ImageNet-R ($\Delta{=}+0.18$,
$p{=}0.38$).

\paragraph{What remains of the Fisher weighting.} One sentence in Sec.~5.2
and one in the conclusion: the rank-$k$ basis protects the measured important directions and
leaves the orthogonal complement free, and it costs 182.8\,MiB of retained state where the
isotropic penalty costs one scalar per layer. We say plainly that we could not show this
matters empirically, and that a practitioner who wants only the accuracy should prefer the
simpler penalty.

\change Title; abstract; highlights; Table~1 (isotropic row);
Sec.~5.2; conclusion.

\subsection*{R8-4 --- The learned-head result: explanation or bias correction}

\begin{rquote}
The new learned-head results (Table 2) show only 7.2--9.2\% accuracy on CIFAR-100 for every
method, which is far below what a linear head on ViT-B/16 features normally achieves and
suggests a classifier failure (for example task-recency bias) rather than a budget effect.
Since this result is now used to justify the NCM protocol, please consider explaining it or
reporting a head with bias correction.
\end{rquote}

\resp You identified the cause exactly. We took the explanation route and report the diagnostic
that establishes it, rather than a bias-corrected head, for a reason given below.

\paragraph{The diagnostic.} We decompose the final model's learned-head accuracy by task
recency tier on the CIFAR-100 sequence, averaged over the same seeds as Table~2:
FOLoRA's head is $97.9\%$ accurate on the last task, $59.9\%$ on the one before it, and
$1.1\%$ averaged over the first eighteen. On ImageNet-R the same three tiers are
$91.1\%$, $43.0\%$ and $0.5\%$. The bulk of the sequence therefore sits at the $1\%$ chance
level ($0.5\%$ on ImageNet-R): a monotone recency gradient that ends at chance, which is the
signature of a classifier
that has moved its decision boundary onto the most recent classes, not of a head that
underfits the representation: the same features reach $77.87\%$ under NCM.

\paragraph{It is uniform across methods.} Every method in Table~2 has the same three-tier
shape: the last-task tier is $97.5$--$97.9\%$ on CIFAR-100 and $85.7$--$91.1\%$ on
ImageNet-R, the remaining eighteen lie between $0.5$ and $1.1\%$ on CIFAR-100 and $0.3$ and
$1.0\%$ on ImageNet-R, and the second-last tier --- the one that varies most between methods,
$37.4$--$69.6\%$ on CIFAR-100 --- is still far above both. The
$7.2$--$9.2\%$ total therefore does not rank the representations and is not evidence for or
against any method; it is evidence about the head.

\paragraph{Why not bias-correct.} A bias correction that is tuned per method (for example a
method-specific temperature or a per-task logit offset) would reintroduce exactly the
method-dependent decision rule that the common NCM protocol removes, and the comparison would
then measure the correction as much as the representation. A single correction shared by all
rows would be defensible, but it would give a second, weaker decision rule whose ordering
would have to be reported in full; the paper already has one such audit (the RanPAC
re-scoring of Sec.~5.3), which corrects the decision rule without training
anything. The learned head is therefore reported as a diagnostic and not used for any
conclusion.

\change Sec.~5.1 (metrics paragraph, recency tiers, and the
cross-reference to the learned-head table).

\subsection*{R8-5 --- ``Replay-free''}

\begin{rquote}
The revision newly discloses that class prototypes are recomputed with the terminal backbone by
re-reading the training data of all past tasks, yet the abstract still describes FOLoRA as
``replay-free''. Please consider qualifying this term (for example ``replay-free during
training'') in the abstract and highlights.
\end{rquote}

\resp Accepted. The abstract now reads ``replay-free during training'', the introduction uses
the same qualified phrase where the method is first summarised, and the conclusion says ``no
replay during training''. The highlights do not use the term at all; the fifth states the
attribution instead --- that the margin comes from sharing the adapter rather than from the
Fisher weighting. The evaluation-time access is stated where the protocol is defined
(Sec.~5.1), with the measured cost of removing it --- the strictly sequential
prototype protocol --- and in the limitations.

\change Abstract; Sec.~1; Sec.~5.1; conclusion.

\subsection*{R8-6 --- Validation-based re-selection, the TOST margin, and ``pre-registered''}

\begin{rquote}
The revision also newly states that hyperparameters were selected on the test set without a
held-out split, and the TOST margin ($\delta{=}1.0$) was set after the reproducibility floor
was measured. With the margins now within about one point of several alternatives, please
consider at least one validation-based re-selection on one benchmark, justifying the margin
independently of the observed data, and removing the term ``pre-registered'' unless a dated
record can be provided.
\end{rquote}

\resp Three separate things; we did all three.

\paragraph{(a) Validation-based re-selection.} We re-ran the selection on CIFAR-100 with a
held-out split that no candidate configuration ever trains on:
\begin{itemize}
\item For each task and each class, $10\%$ of that class's \emph{training} samples are held
  out --- 50 of 500 per class, 5{,}000 samples in total. The holdout indices are derived
  deterministically from (seed, task, class), so all candidate $\lambda$ at a given seed
  share the \emph{same} holdout; changing $\lambda$ changes what the model trains on, never
  what it is selected on.
\item We train $\lambda\in\{1,3,10\}$ on seeds $0,1,2$ with every other hyperparameter fixed
  at the main configuration, and evaluate the resulting checkpoint under the same NCM rule
  used in the main table, on the held-out samples only. $\lambda$ is then chosen by held-out
  accuracy, and we report that $\lambda$'s accuracy on the official test set. The test set is
  not consulted at any point in the selection.
\item The selection does not differ: held-out accuracy picks $\lambda{=}3$ over
  $\lambda{=}1$ and $\lambda{=}10$ ($76.70\pm0.40$ against $74.73\pm0.86$ and
  $73.20\pm1.90$), on the mean over seeds and at each of the three seeds individually. The
  ordering of the three candidates on the holdout ($3>1>10$) is also their ordering on the
  test set --- the property this check was meant to establish. The selected configuration's
  test accuracy under this protocol is $77.04\pm0.31\%$, against $78.37\pm0.84\%$ for the same
  $\lambda$ in Table~1; the two are not comparable, because these models train on $90\%$ of
  the training data, and we report the new number as the answer to your question rather than
  as a second estimate of the method.
\end{itemize}
The script, the per-run JSON and the selection record are in the replication package
(\texttt{scripts/val\_select.py}, \texttt{reports/val\_selection.json}). This is, we
acknowledge, a re-selection we did \emph{after} the fact --- it answers whether the protocol
selects the same configuration, not whether we followed it from the start. Sec.~5.1
now states the protocol the main configuration came from, including that the grid was run on
test accuracy, and points to this experiment as the validation-based check.

\paragraph{(b) The TOST margin, justified independently of the data.} The margin is now fixed
on a substantive basis: one point of average accuracy on a 100-class benchmark is the smallest
difference we would treat as of practical consequence for a method choice, and $1.0$ is the
round value at that scale. The reproducibility floor no longer \emph{defines} the margin; it
enters only as a constraint the margin must clear --- a margin set at or below the $0.89$-point
run-to-run difference would claim an equivalence no number of seeds could establish, so the
margin must be above it, and $1.0$ is the smallest round value above it. We also report the
sensitivity: at $\delta{=}0.75$ the tuned EWC-LoRA points are already equivalent on both
benchmarks ($p=0.032$--$0.055$), at $\delta{=}1.25$ they remain equivalent, and at
$\delta{=}0.5$ none is equivalent on accuracy --- the expected outcome when the margin is
narrower than the measured floor, and the reason that value is not used. The headline verdict does not depend on the
choice within the range the floor permits.

\paragraph{(c) ``Pre-registered''.} You are right that the term was used where no record
exists. We audited every occurrence. Where a dated record \emph{does} exist, it is the
project decision log, entry dated \textbf{23 September 2026, 19:45}, written when 13 of the 23
$\lambda$-grid runs had finished and seeds 3--4 had not yet landed. It fixes, before those
runs: the main configuration ($k{=}64$, $\lambda{=}3$); the comparison point for the headline
verdict (EWC-LoRA at its own default $\lambda{=}100$, paired $t$-test on common seeds); and
the wording thresholds for that verdict. Quoted verbatim (the log is in Chinese; translation
ours):

\begin{CJK}{UTF8}{gbsn}
\begin{quote}\small
\textbf{13.3 冻结判决规则（在 seed 3/4 落地之前写死，防 garden-of-forking-paths）}\\
剩下的 seed 3/4 落地后，\textbf{按下面这条规则判决，不再改口径}：\\
1. \textbf{论文一律报告完整 $\lambda$ 曲线}（$\lambda\in\{1,3,10\}$ 各 5 seeds 的
mean$\pm$std + 误差棒），\textbf{不隐藏 $\lambda=10$ 不显著这件事} \ldots\\
2. 主配置 = \textbf{k=64, $\lambda$=3}（判决前已定，不是看到结果才定）。\\
3. 检验用\textbf{配对 t（共同 seed）}，与 EWC $\lambda$=100（n=10）比：
若 p$<$0.05 且 5 seeds 逐 seed 差全为正 $\rightarrow$ 可写 ``outperforms EWC-LoRA''，并同时报效应量（约 +1.1 点）；
若 0.05$\le$p$<$0.10 $\rightarrow$ 只写 ``matches or slightly exceeds''，\textbf{不用 outperforms}；
若 p$\ge$0.10 $\rightarrow$ 回到「平价」定位 \ldots\\
4. \textbf{绝不做的事}：不因为 $\lambda$=10 某几个 seed 高就换成 $\lambda$=10 \ldots
\end{quote}
\end{CJK}

We do not claim that this record extends further than it does: it does \emph{not} fix the TOST
margin, it does not cover the second benchmark, and the final verdict it produced was
``parity'', not a win. The paper now says this explicitly in Sec.~5.1, and every
other use of ``pre-registered'' was removed --- in particular the claim that the baseline's
$\lambda{=}100$ was a pre-registered comparison point (it is the baseline's default, which is
why it is the fair comparison point, and the text now says that instead), and the phrase
``before any seed was run'', which was not true of the log entry. The full log entry, the
document's timestamp and its position in the file are in the replication package.

\change New Sec.~5.1 paragraph on how the main configuration was selected and on
what the dated record covers; TOST paragraph rewritten; Table~1 caption
(``judging baseline $\lambda{=}100$'' now ``the baseline's own default $\lambda{=}100$'');
new validation-selection experiment (Sec.~5.1,
Table~A.11, \texttt{reports/val\_selection.json}).

\subsection*{R8-7 --- Unrelated metaheuristic references}

\begin{rquote}
In response to another reviewer, five metaheuristic optimization papers whose topics (5G QoS
classification, desalination modelling, energy forecasting and new optimizers) are unrelated to
continual learning were added in a single sentence (Sec.~5.1, references [76]--[80]). Please
consider keeping only references that directly support the text.
\end{rquote}

\resp Accepted. The five references are removed from the manuscript and from the bibliography,
and the sentence that carried them now makes the same point without citation: we enumerate the
grid rather than search it with a derivative-free optimizer, because a search whose trajectory
cannot be enumerated would make the baseline tuning harder to audit. These five were added in
the previous round at another reviewer's suggestion; re-reading them against your criterion, we
agree that none of them supports the sentence, and we removed the whole cluster rather than keep
part of it.

The rest of the bibliography was re-audited under the same criterion. It now has 78 entries,
every one cited in the text, with no uncited entry and none cited only in the appendix. Three
recent works that could look equally peripheral were kept after that check, because each is
cited for a specific claim it supports: the LoRA/QLoRA/full-fine-tuning comparison and
TabEng-QLoRA (Related Work, where the deployment-budget accounting is introduced) and SQ-LoRA
(Sec.~5.5, next to the $k{=}16$ recommendation).

\change Sec.~5.1 (configuration-selection paragraph); bibliography (five entries
removed, 78 entries re-audited: all cited in the text); reference numbering throughout.
